% Book-local P03.1 upper-regularization defect localization proof,
% 19 September 2026.  No novelty or priority claim is made.

\begin{proof}[\Cref{thm:p031-usc-defect-localization}]
By construction $u=s^*\geq s$.  An ordinary open set $U\subseteq X$
satisfies $u=s$ on $U$ exactly when $U\cap D(s)=\varnothing$.
It follows that the union of all such open sets is
\[
                         X\setminus\overline{D(s)}=G(s),
\]
which proves the maximality assertion.  Since $u$ is upper
semicontinuous, the equality $s=u$ on $G(s)$ also shows that
$s|_{G(s)}$ is upper semicontinuous.

Let $\mathcal F(I)$ be the directed set of nonempty finite subsets of
$I$, ordered by inclusion.  For $F\in\mathcal F(I)$, put
\[
                            m_F:=\max_{i\in F}\varphi_i.
\]
The cone property
\cite[Proposition~4.2(3)]{GR25Classical} makes every $m_F$ classically
psh.  On $G(s)$, the net $(m_F)_F$ converges pointwise to
\[
                         s|_{G(s)}=u|_{G(s)}.
\]
Its limit is upper semicontinuous, so the pointwise-limit theorem
\cite[Theorem~6.5]{GR25Classical} proves that $u|_{G(s)}$ is classically
psh.

Finally, $s$ is upper semicontinuous if and only if $s=s^*=u$, which is
equivalent to $D(s)=\varnothing$.  The preceding conclusion then applies
with $G(s)=X$.
\end{proof}

\begin{proof}[\Cref{cor:p031-completed-algebraic-closure-defect-support}]
Equation \eqref{eq:p031-no-completed-algebraic-closure-defect} gives
\[
                     u_L=u\circ\pi_{L/K},
 \qquad
                     s_L=s\circ\pi_{L/K}.
\]
The strict inequality defining the defect set therefore gives
$D_L=\pi_{L/K}^{-1}(D_K)$.

Continuity of $\pi_{L/K}$ gives
\[
 \overline{\pi_{L/K}^{-1}(D_K)}
       \subseteq\pi_{L/K}^{-1}(\overline{D_K}).
\]
For the reverse inclusion, fix $y$ over a point of
$\overline{D_K}$.  Every ordinary open neighborhood $N$ of $y$ has
open image by \Cref{thm:p031-completed-algebraic-closure-reduction}; this
image meets $D_K$, so $N$ meets $\pi_{L/K}^{-1}(D_K)$.  Hence
$y\in\overline{D_L}$, which proves the second equality.  Taking
complements gives the final identity.
\end{proof}
